% !LW recipe = latexmk (XeLaTeX, automatic references)
% !TeX encoding = UTF-8
% 编译方法：latexmk -xelatex -interaction=nonstopmode -halt-on-error -outdir=build 第二问模型.tex
\documentclass[UTF8,a4paper,zihao=-4,fontset=none]{ctexart}

% 字体采用三级自动选择机制，使模板能够在不同电脑上直接编译。
% 先检查 Windows 标准字体目录；再检查 macOS 版 Microsoft Word 的字体；均不可用时回退到 TeX Live 自带的 Fandol 字体。
% 无论选择哪一组字体，主字体的 BoldFont 都显式绑定到对应黑体，确保 \textbf 下的中文字体正确。
\newcommand{\UseWindowsMicrosoftFonts}{
  \setCJKmainfont{simsun.ttc}[Path={C:/Windows/Fonts/},BoldFont=simhei.ttf]
  \setCJKsansfont{simhei.ttf}[Path={C:/Windows/Fonts/},BoldFont=simhei.ttf]
  \setCJKmonofont{simsun.ttc}[Path={C:/Windows/Fonts/}]
  \setCJKfamilyfont{paperhei}{simhei.ttf}[Path={C:/Windows/Fonts/},BoldFont=simhei.ttf]
  \PackageInfo{paper-template}{Using Windows SimSun and SimHei}
}

\newcommand{\UseOfficeMicrosoftFonts}{
  \setCJKmainfont{Simsun.ttc}[
    Path={/Applications/Microsoft Word.app/Contents/Resources/DFonts/},
    BoldFont=SimHei.ttf
  ]
  \setCJKsansfont{SimHei.ttf}[
    Path={/Applications/Microsoft Word.app/Contents/Resources/DFonts/},
    BoldFont=SimHei.ttf
  ]
  \setCJKmonofont{Simsun.ttc}[
    Path={/Applications/Microsoft Word.app/Contents/Resources/DFonts/}
  ]
  \setCJKfamilyfont{paperhei}{SimHei.ttf}[
    Path={/Applications/Microsoft Word.app/Contents/Resources/DFonts/},
    BoldFont=SimHei.ttf
  ]
  \PackageInfo{paper-template}{Using Microsoft Office SimSun and SimHei}
}

\newcommand{\UsePortableFandolFonts}{
  \setCJKmainfont{FandolSong-Regular.otf}[cmap=UniGB-UTF16-H,BoldFont=FandolHei-Regular.otf]
  \setCJKsansfont{FandolHei-Regular.otf}[cmap=UniGB-UTF16-H,BoldFont=FandolHei-Regular.otf]
  \setCJKmonofont{FandolFang-Regular.otf}[cmap=UniGB-UTF16-H]
  \setCJKfamilyfont{paperhei}{FandolHei-Regular.otf}[cmap=UniGB-UTF16-H,BoldFont=FandolHei-Regular.otf]
  \PackageWarningNoLine{paper-template}{SimSun and SimHei not found; using portable Fandol fonts}
}

\newcommand{\UseOfficeOrPortableFonts}{
  \IfFileExists{/Applications/Microsoft Word.app/Contents/Resources/DFonts/Simsun.ttc}{
    \IfFileExists{/Applications/Microsoft Word.app/Contents/Resources/DFonts/SimHei.ttf}{
      \UseOfficeMicrosoftFonts
    }{
      \UsePortableFandolFonts
    }
  }{
    \UsePortableFandolFonts
  }
}

% 将下面的 0 改为 1，可强制使用完全随 TeX Live 提供的 Fandol 字体进行兼容性测试。
\providecommand{\ForcePortableFonts}{0}
\ifnum\ForcePortableFonts=1
  \UsePortableFandolFonts
\else
  \IfFileExists{C:/Windows/Fonts/simsun.ttc}{
    \IfFileExists{C:/Windows/Fonts/simhei.ttf}{
      \UseWindowsMicrosoftFonts
    }{
      \UseOfficeOrPortableFonts
    }
  }{
    \UseOfficeOrPortableFonts
  }
\fi
\newcommand{\heifont}{\CJKfamily{paperhei}}

% 页面、行距与段落。
\usepackage[a4paper,top=2cm,bottom=2cm,left=2cm,right=2cm,footskip=1.2cm]{geometry}
\usepackage{setspace}
\singlespacing
\setlength{\parindent}{2em}
\setlength{\parskip}{0pt}
\setlength{\emergencystretch}{3em}
\usepackage{indentfirst}

% 数学公式、表格与插图。
\usepackage{amsmath,amssymb,mathtools,bm}
\usepackage{graphicx}
\graphicspath{{./}}
\usepackage{booktabs,tabularx,array,multirow}
\usepackage[table]{xcolor}
\usepackage{float}
\usepackage{caption}
\usepackage{subcaption}
\usepackage{enumitem}

% 标题、页码与交叉引用。
\usepackage{titlesec}
\usepackage{fancyhdr}
\usepackage[hidelinks]{hyperref}
\usepackage{cleveref}

% 一级标题使用中文序号并居中，下级标题使用阿拉伯数字层级编号并左对齐。
\renewcommand{\thesection}{\arabic{section}}
\renewcommand{\thesubsection}{\arabic{section}.\arabic{subsection}}
\renewcommand{\thesubsubsection}{\arabic{section}.\arabic{subsection}.\arabic{subsubsection}}
\setcounter{secnumdepth}{3}
\titleformat{\section}[block]
  {\centering\heifont\bfseries\zihao{4}}
  {\chinese{section}、}{0pt}{}
\titleformat{\subsection}[hang]
  {\raggedright\heifont\bfseries\zihao{-4}}
  {\thesubsection}{0.8em}{}
\titleformat{\subsubsection}[hang]
  {\raggedright\heifont\bfseries\zihao{-4}}
  {\thesubsubsection}{0.8em}{}
\titlespacing*{\section}{0pt}{2.8ex plus 0.5ex minus 0.2ex}{1.7ex}
\titlespacing*{\subsection}{0pt}{1.6ex plus 0.3ex minus 0.1ex}{0.8ex}
\titlespacing*{\subsubsection}{0pt}{1.4ex plus 0.2ex minus 0.1ex}{0.7ex}

% 页码从摘要页开始，位于页脚中央。
\pagestyle{fancy}
\fancyhf{}
\fancyfoot[C]{\zihao{5}\thepage}
\renewcommand{\headrulewidth}{0pt}
\renewcommand{\footrulewidth}{0pt}

% 图表及列表的常用设置。
\captionsetup{font=small,labelfont=bf,labelsep=quad}
\setlist[itemize]{itemsep=0.35\baselineskip,topsep=0.8ex,parsep=0pt,partopsep=0pt}
\newcolumntype{C}[1]{>{\centering\arraybackslash}p{#1}}
\newcolumntype{L}[1]{>{\raggedright\arraybackslash}p{#1}}
\newcolumntype{Y}{>{\raggedright\arraybackslash}X}
\newcolumntype{Z}{>{\centering\arraybackslash}X}

% 请优先修改下面三个字段。
\newcommand{\papertitle}{第二检测点选择：统一损失下的期望与最坏准则}
\newcommand{\paperabstract}{}
\newcommand{\paperkeywords}{}

\hypersetup{
  pdftitle={\papertitle},
  pdfauthor={}
}

\begin{document}
\begin{center}
{\heifont\bfseries\zihao{3}\papertitle\par}
\end{center}

\section*{记号约定}

本文件独立呈现论文问题二的定稿模型。设目标圆域为 \(\Omega=\{g\in\mathbb R^2:\|g\|\le1800\}\)，\(g\) 为固定未知的全向干扰源位置，\(s_i\) 为检测点，\(\theta_i\) 为从正东方向逆时针度量的正常示向度。角度以弧度计，距离以米计；测角误差界为 \(\alpha=\pi/180\)，强信号距离阈值为 \(r_{\mathrm{strong}}=5\)，光学定位距离为 \(r_{\mathrm{opt}}=20\)，固定未知的接收半径为 \(R\in[1000,1500]\)。

正常示向度对应的正向角域记为 \(W_i=\{s_i+t(\cos(\theta_i+\delta),\sin(\theta_i+\delta)):t\ge0,\ |\delta|\le\alpha\}\)，\(\mathbf1_A\) 表示集合 \(A\) 的指示函数。正文与论文的问题二章节保持一致，所需其余记号随推导定义。

\section{问题二的模型与求解}

该问题的任务是基于第一个检测点获得的信息制定第二个检测点的选择策略。由于选点时并不确定第二次会收到什么反馈，我们需要先比较各个检测点可能带来的结果。针对本问题的目标，我们在此仅以定位精确度来衡量定位效果，而暂不考虑执行策略消耗时间的差异。

据此，为了比较不同检测结果，本文用“定位损失”衡量检测后还剩下多少位置不确定性。若尚不能在第二检测点原地完成光学精确定位，就以剩余可行区域的最小包围圆半径作为定位损失：半径越小，说明源的位置被限制在越小的范围内，定位效果越好。若整个剩余区域都已位于当前点$20$米以内，则能够原地完成光学精确定位，定位损失取零。这里的损失是根据可行区域定义的评价指标，而不是指估计位置与真实位置之间的误差。

进一步，本文设计了两种评价方法。\textbf{基于贝叶斯更新的期望定位损失模型}先预测不同反馈出现的概率，再计算平均定位损失，选择平均损失最小的位置。\textbf{最小化最坏剩余半径模型}则先找出每个位置可能遇到的最不利反馈，再选择最坏损失最小的位置。两种方法使用相同的定位区域和损失定义，分别回答“平均而言在哪里检测更好”和“在哪里检测能获得更好的最坏情形保证”。

\subsection{准备工作：由观测信息构造定位区域}

首先设第二个检测点为 \(q\)，它的位置是本问需要决定的量。为简化坐标表示，利用圆的旋转对称性，不妨将首次检测点写成 \(s_1=(a,0)\)，首次示向度记为 \(\beta\)，其中 \(a\ge0\)、\(\beta\in[-\pi,\pi)\) 均为确定的参数。若在 \(q\) 处得到正常示向度 \(\theta\)，则令 \(s_2=q\)、\(\theta_1=\beta\)、\(\theta_2=\theta\)，并沿用问题一的角域记号 \(W_1,W_2\)。源位置 \(g\) 到两次检测点的距离分别简记为 \(d_1(g)=\|g-s_1\|\)、\(d(g;q)=\|g-q\|\)。

首次测得正常示向度，说明源位于对应角域内，能够被接收，并且尚未近到触发强信号。因此，首次观测后的源位置可行区域为 \(K_1=\{g\in\Omega\cap W_1:r_{\mathrm{strong}}<d_1(g)\le1500\}\)。对于其中一个可能位置 \(g\)，接收半径还必须满足 \(\max\{1000,d_1(g)\}\le R\le1500\)。

在 \(q\) 处进行第二次检测，可能得到强信号（源距$q$不超过5米）、无信号（源距大于接收半径）或正常示向度，分别记为 \(\mathrm H\)、\(\mathrm N\) 和 \(\theta\)。
需要注意，无信号也能提供信息。由于首次探测已成功接收，且两次检测对应同一个固定接收半径，第二次无信号说明源距 \(q\) 比源距 \(s_1\) 更远，同时超过1000米。将这些新信息与 \(K_1\) 取交集，得到三种反馈后的剩余区域

\begin{equation}
\begin{aligned}
K_{\mathrm H}(q)&=\{g\in K_1:d(g;q)\le r_{\mathrm{strong}}\},\\
K_{\mathrm N}(q)&=\{g\in K_1:d(g;q)>\max\{1000,d_1(g)\}\},\\
K_\theta(q)&=\{g\in K_1\cap W_2:r_{\mathrm{strong}}<d(g;q)\le1500\}.
\end{aligned}
\label{eq:q2-pair-regions}
\end{equation}

其中，\(W_2\) 由检测点 \(q\) 和读数 \(\theta\) 确定。若某种反馈对应的区域为空，它就不可能在已有观测条件下出现，后续不计入评价。将其余可能出现的反馈记为集合 \(F(q)\)。

沿用问题一的方法，用最小包围圆半径 \(\rho(K)\) 衡量区域 \(K\) 还有多大，半径越小，源的位置就确定得越精细。第二次反馈为 \(z\) 时，将剩余区域 \(K_z(q)\) 的这一半径简记为 \(\rho_z(q)\)。此外，用 \(d_{\max}(q;K)\) 表示 \(q\) 到区域内最远位置的距离，用它判断能否在 \(q\) 原地完成光学定位。

若第二个检测点距 \(K_1\) 超过1500米，无论源在区域内何处，都只能得到无信号，因而无法缩小已有区域。因此，将距 \(K_1\) 不超过1500米、且不同于首次检测点的位置作为搜索范围，记为 \(Q\)。这里限制的是检测点到可能源位置的距离，检测点本身可以位于目标圆域 \(\Omega\) 外。

\subsection{反馈后的统一定位损失}

评价一次检测是否有用，要看得到反馈后还剩下多少位置不确定性，这就是定位损失。具体地，定位损失以米为单位，数值越小，表示检测后的定位效果越好。根据第二次检测的反馈，计算规则如下。

\begin{enumerate}[label=(\arabic*),leftmargin=2.2em]
  \item \textbf{强信号。} 源距当前检测点不超过5米，已在20米的光学定位范围内，因此可以原地确定源的准确位置，取 \(L_{\mathrm H}(q)=0\)。
  \item \textbf{无信号。} 源距当前检测点超过1000米，无法原地完成光学定位。因此，以无信号反馈后剩余区域的最小包围圆半径作为损失，取 \(L_{\mathrm N}(q)=\rho_{\mathrm N}(q)\)。
  \item \textbf{正常示向度。} 根据新读数 \(\theta\) 更新定位区域，再判断该区域是否整体位于当前检测点的20米范围内。若是，则可以原地完成光学定位，取 \(L_\theta(q)=0\)；否则，以更新后区域的最小包围圆半径作为损失，取 \(L_\theta(q)=\rho_\theta(q)\)。
\end{enumerate}

\subsection{基于贝叶斯更新的期望定位损失模型}

对于同一个检测点，有的反馈会使干扰源候选区域明显缩小，有的反馈则可能帮助不大。本模型先计算各种反馈有多大可能出现，再按这些概率求平均损失。为此，需要先用首次观测更新对源位置和接收半径的判断，再预测第二次检测的结果。

\subsubsection{固定未知量的贝叶斯描述}

源位置、接收半径和各检测地点的误差都有固定取值，但我们事先不知道这些取值。\textbf{贝叶斯模型用概率分布描述我们对固定未知量的认识。} 新观测到来后，改变的是对其可能取值的判断，未知量本身并不随之改变。因此，为地点误差设置概率分布，与题目规定的“同一地点重复检测不改变误差”并不矛盾。

按照前文模型假设，源位置 \(G\) 在 \(\Omega\) 内按面积均匀分布，接收半径 \(R\sim\operatorname{U}[1000,1500]\)，两个不同检测点的误差 \(E_1,E_2\sim\operatorname{U}[-\alpha,\alpha]\)，观测前认为这些未知量相互独立。均匀性和独立性是为了计算反馈概率而增加的假设，不能仅由题目给出的误差界与接收范围推出。该方法得到的平均效果也以这些假设为前提。

同一地点的多次检测始终共享同一个误差，不能重新抽样，也不能靠重复测量取平均消除误差。本模型只预测不同于首次检测点的新位置。对于观测后产生的关联，例如较远的源位置必须对应较大的接收半径，后续计算会一并保留。

\subsubsection{首次观测后的概率更新}

将首次正常示向度信息记为 \(I_1=(s_1,\beta)\)。可行区域 \(K_1\) 告诉我们哪些位置没有被排除，但其中的位置未必同样可能。例如，源距不超过1000米时，题目允许的任何接收半径都能解释首次接收；距离超过1000米后，只有较大的接收半径才能解释这一观测。因此，应根据接收条件调整不同位置的概率。

贝叶斯更新将原有概率与观测的似然相乘，再归一化。似然表示假定某个源位置成立时，得到已有观测的可能程度。在 \(K_1\) 内，均匀测角误差对应的角度似然相同，位置权重的差别来自能否接收到信号。距离为 \(d\) 时，这一接收概率为

\begin{equation}
 w(d)=\Pr(R\ge d)=
 \begin{cases}
 1,&d\le1000,\\
 (1500-d)/500,&1000<d<1500,\\
 0,&d\ge1500.
 \end{cases}
\label{eq:q2-pair-weight}
\end{equation}

将上述权重用于 \(K_1\) 内的各个位置，并使总概率等于1，得到首次观测后的源位置密度

\begin{equation}
 p_1(g)=\frac{\mathbf1_{K_1}(g)w(d_1(g))}{C_1},
 \qquad C_1=\int_{K_1}w(d_1(g))\,\mathrm dg>0.
\label{eq:q2-pair-posterior}
\end{equation}

其中 \(C_1\) 是归一化系数。该结果说明，距离超过1000米后，越远的位置得到的权重越低，因为能解释首次接收的半径取值越来越少。

接收半径本身也需要更新。若源位于 \(g\)，首次接收已经排除了小于 \(d_1(g)\) 的半径。因此，对于后验密度为正的位置，有 \(R\mid G=g,I_1\sim\operatorname{U}[\max\{1000,d_1(g)\},1500]\)。预测第二次检测时，要使用这个更新后的范围，不能再次把 \(R\) 当作与源位置无关的均匀量。

\subsubsection{第二次反馈的概率预测}

现在固定一个待评价的检测点 \(q\)，考虑源的每个可能位置 \(g\)。因为接收半径固定，两次都能接收要求半径至少达到两次源距中较大的一个。记这个距离为 \(d_\vee(g;q)=\max\{d_1(g),d(g;q)\}\)，则两次均能接收的概率权重为 \(w(d_\vee)\)。从首次能够接收的权重 \(w(d_1)\) 中减去它，就得到首次接收、第二次无信号的权重。

再结合三种反馈各自的位置约束，得到用于计算反馈概率的权重函数

\begin{equation}
\begin{aligned}
 b_{\mathrm H}(g;q)&=\mathbf1_{K_{\mathrm H}(q)}(g)w(d_1(g)),\\
 b_{\mathrm N}(g;q)&=\mathbf1_{K_1}(g)\bigl[w(d_1(g))-w(d_\vee(g;q))\bigr],\\
 b_\theta(g;q)&=\frac{\mathbf1_{K_\theta(q)}(g)}{2\alpha}w(d_\vee(g;q)).
\end{aligned}
\label{eq:q2-pair-feedback-weights}
\end{equation}

其中，正常示向度的读数还受新地点测角误差影响。按照不同地点误差独立且均匀的假设，允许的角度范围内密度为 \(1/(2\alpha)\)，因此第三项包含这个因子。将各个可能源位置的权重积分，并除以首次观测的归一化系数，得到

\begin{equation}
\begin{aligned}
 P_{\mathrm H}(q)&=\frac{1}{C_1}\int_{K_1}b_{\mathrm H}(g;q)\,\mathrm dg,
 &P_{\mathrm N}(q)&=\frac{1}{C_1}\int_{K_1}b_{\mathrm N}(g;q)\,\mathrm dg,\\
 f_\Theta(\theta;q)&=\frac{1}{C_1}\int_{K_1}b_\theta(g;q)\,\mathrm dg.
\end{aligned}
\label{eq:q2-pair-prediction}
\end{equation}

\(P_{\mathrm H}(q)\) 和 \(P_{\mathrm N}(q)\) 分别是强信号与无信号出现的概率。正常示向度可以连续变化，因此用密度 \(f_\Theta(\theta;q)\) 描述；将它在某个角度区间内积分，就得到读数落入该区间的概率。在全部角度上积分所得的正常反馈概率，与前两种概率相加等于1。

\subsubsection{最小化期望定位损失的选点决策}

至此，对于每个待选位置 \(q\)，我们既知道不同反馈会留下多大的定位损失，也知道这些反馈出现的概率。以 \(Z\) 表示尚未得到的第二次反馈，将损失按预测概率求平均，记为 \(J_{\mathrm E}(q)\)。第二检测点的决策就是在搜索范围 \(Q\) 内，使这个平均损失尽可能小，即

\begin{equation}
 \min_{q\in Q}\quad J_{\mathrm E}(q)=\mathbb E[L_Z(q)\mid I_1,q]
 =P_{\mathrm N}(q)\rho_{\mathrm N}(q)
 +\int_0^{2\pi}f_\Theta(\theta;q)L_\theta(q)\,\mathrm d\theta.
\label{eq:q2-pair-expected-objective}
\end{equation}

右端第一项对应无信号的结果，第二项将各种正常示向度对应的损失按概率密度积分。强信号能够直接完成原地定位，损失为零，所以不再单列一项。求解这个最小化问题，就得到所设先验下平均定位效果最好的检测位置；但平均效果好，并不保证每一次反馈都能带来较小的损失。

\subsection{最小化最坏剩余半径的模型}

如果希望避免某些反馈带来很差的定位结果，可以改用最坏情形准则。对于每个检测点 \(q\)，先考察所有可能反馈，将其中最大的定位损失记为 \(J_{\mathrm W}(q)\)。再比较不同检测点，选择这个最大损失最小的位置，即

\begin{equation}
 \min_{q\in Q}\quad J_{\mathrm W}(q)=\sup_{z\in F(q)}L_z(q).
\label{eq:q2-pair-worst-objective}
\end{equation}

这个模型不需要假设各反馈出现的概率。只要某种反馈符合题意、没有被已有观测排除，就必须考虑其后果，而不能因为它很少出现就忽略它。因此，该方法给出的是不利情况下仍能达到的定位效果。

对于同一个检测点，平均损失不会超过最大的可能损失，即 \(J_{\mathrm E}(q)\le J_{\mathrm W}(q)\)。但两者最小值所在的位置可能不同：一个位置可能平均效果更好，另一个位置则可能在最不利反馈下表现更好。

\subsection{候选区域与实际选点策略}

\subsubsection{从最优检测点扩展到候选区域}

题目还要求给出候选区域，而不只是一个最优点。为此，在求解上述最小化问题后，保留所有定位效果与最优值相差不大的位置。用 \(J\) 统一表示 \(J_{\mathrm E}\) 或 \(J_{\mathrm W}\)，最优损失记为 \(J^*=\inf_{q\in Q}J(q)\)。若允许比最优值多出 \(\epsilon\) 米的损失，则候选区域为

\begin{equation}
 Q_{J,\epsilon}=\{q\in Q:J(q)\le J^*+\epsilon\}.
 \label{eq:q2-candidate-region}
\end{equation}

例如，若最小期望损失为8米，允许增加1米，就保留所有期望损失不超过9米的位置。这个例子仅用于说明区域的定义，不是本题的计算结果。区域由各个位置的定位效果决定，可能分成几片；这些部分都应保留，不能将中间不满足要求的位置也连接进来。

两种模型分别计算自己的最优值和候选区域。期望模型的区域提供平均效果接近最优的选择，最坏模型的区域提供最坏损失接近最优的选择。绘图时将二者叠加，便于比较；不要求实际检测点一定在交集中，因为交集可能为空。

\subsubsection{容差选择与区域计算}

容差 \(\epsilon\) 表示愿意让定位损失增加多少米，以获得更多可选位置。本文暂以1米为默认值，并用0.5米、2米进行对照。它是可以调整的策略参数，需要通过后续计算比较区域大小和移动距离后判断是否合适。采用绝对容差，也使最优损失为零时仍能保留附近效果相近的位置。

数值求解时，先在 \(Q\) 内设置网格并计算各点的 \(J\)，找到当前最低损失，再按容差筛选候选位置。对损失较低的区域和接近筛选边界的位置加密网格，并提高反馈概率积分或最坏反馈搜索的精度。不能因为网格的几个顶点合格，就认定整个网格都合格；边界附近尚不能确认的位置应继续检查，最终推荐点也要单独复算。

结果图应标出首次检测点、源位置区域 \(K_1\)、两种候选区域及推荐点，同时注明损失基准和容差。各区域随首次观测 \((a,\beta)\) 重新计算。若只完成了有限精度的搜索，就将结果称为数值近似候选区域；网格加密后结果趋于稳定，可以支持数值结果的可信度，但本身不能证明已经找到连续空间中的全局最优值。

\subsubsection{从候选区域中确定实际检测点}

前两种模型首先解决定位效果问题。得到候选区域后，其中的位置都满足允许的损失要求，可以再考虑移动是否方便。本文优先选择离首次检测点最近的位置；在最小值能够取到时，实际推荐点满足

\begin{equation}
 q^*\in\arg\min_{q\in Q_{J,\epsilon}}\|q-s_1\|.
 \label{eq:q2-candidate-selection}
\end{equation}

这一步最小化的是移动距离，而前面最小化 \(J\) 是为了确定定位效果的基准，二者作用不同。实际推荐点不一定取得最小定位损失，但其损失仍在设定容差以内。若多个点距离相同，则优先选择定位损失较小者；计算时只从已经复核的候选点中选择。


\end{document}
